\chapter{Restricted duals and differential jets}
\label{cyjets:ch}
\begingroup
\def\eps{\epsilon}

For a polynomial $F(x)$, evaluation of its $m$th derivative at zero
is a coefficient functional. Transposing differentiation sends this
functional to evaluation of the next derivative. Thus a locally nilpotent
operation on polynomials becomes an infinite translation orbit on their
restricted dual. This elementary change of carrier permits a singular
field to arise from a spectral Hall cooperation.

We first construct that field for the three-loop quiver with zero
potential. Its contraction is the full inverse Euler factor, including
an even factor invisible to the exchange ratio. The comparison is an
identity on polynomial tests, with explicit factorials, translation sign
and Laurent bounds. The critical Hall coefficients of the preceding
chapter require a further geometric map.

An infinite current orbit also determines the coefficient category for
geometric duality. Finite perfect differential modules admit an
involutive dual, and dualizing a filtered presentation retains an arbitrary
coefficient as a pro-object. Normal support imposes a further condition:
it survives additive convolution but can be lost by an ordinary realized
dual. After proving these statements, we transpose the coproduct of a
full bar construction. The resulting associative algebra has a transported
tensor; its realization by independently defined fields is a separate
comparison.

\section{Derivative evaluation}
\label{cyjets:sec:pairing}

Let $k$ be a field of characteristic zero. On $k[x]$ define
\[
 \delta_m(F)=F^{(m)}(0),\qquad m\geq0.
\]
Then $\delta_m(x^n)=m!\,\delta_{mn}$. The vector space spanned by
the $\delta_m$ is the continuous dual of $k[[x]]$, where $k$ is
discrete and $k[[x]]$ has its $x$-adic topology. It is smaller than
the algebraic dual of $k[x]$: a continuous functional depends on
only finitely many coefficients. If $D=-\partial_x$ and $D^\circ$
denotes precomposition by $D$, then
\begin{equation}\label{cyjets:eq:rank-one-transpose}
 D^\circ\delta_m=-\delta_{m+1}.
\end{equation}
Consequently the finite span of $\delta_0,\ldots,\delta_M$ is not
preserved by $-D^\circ$.

For the Hall example take
\[
 k=\Q(a,b),\qquad h=a+b,\qquad
 R_N=k[\eps]/(\eps^N),\qquad N\geq1.
\]
The parameter $\eps$ is nilpotent and is never inverted. In charge
$d\geq0$, put
\[
 A_d=R_N[x_1,\ldots,x_d]^{\mathfrak S_d},\qquad A_0=R_N.
\]
Let $A_d[m]$ denote its homogeneous root-degree-$m$ part. Each such
part is finite free over $R_N$. Its degree completion and restricted
dual are, respectively,
\begin{equation}\label{cyjets:eq:restricted-dual}
 \widehat A_d=\prod_{m\geq0}A_d[m],\qquad
 A_d^\circ=\bigoplus_{m\geq0}
                 \Hom_{R_N}(A_d[m],R_N).
\end{equation}
The direct sum is the continuous dual of the displayed product,
with the finite-degree filtration and discrete target $R_N$.

Introduce the polynomial state space
\[
 W_N=R_N[J_0,J_1,J_2,\ldots],\qquad
 W_{N,d}=\operatorname{span}_{R_N}
       \{J_{m_1}\cdots J_{m_d}:m_i\geq0\}.
\]
Thus $W_N=\bigoplus_dW_{N,d}$ and $W_{N,0}=R_N1$.
Its ordinary polynomial multiplication concatenates the $J$ factors.
Define a derivation $T$ by $TJ_m=J_{m+1}$ and $T1=0$.
All these states have cochain degree zero. Pair states with Hall
polynomials by
\begin{equation}\label{cyjets:eq:pairing}
 \langle J_{m_1}\cdots J_{m_d},F\rangle
 =\left.\partial_{x_1}^{m_1}\cdots\partial_{x_d}^{m_d}F
                                                    \right|_{x=0}.
\end{equation}
Symmetry of $F$ makes this independent of the ordering of the $m_i$.

\begin{proposition}\label{cyjets:prop:pairing}
The pairing identifies $A_d^\circ$ with $W_{N,d}$. Under this
identification,
\[
 T=-D^\circ,\qquad D=-\sum_{i=1}^d\partial_{x_i}.
\]
\end{proposition}
\begin{proof}
For a nondecreasing list $\lambda=(\lambda_1,\ldots,\lambda_d)$,
let $m_\lambda$ be the sum of the distinct monomials obtained by
permuting the exponents of $x_1^{\lambda_1}\cdots x_d^{\lambda_d}$.
These form a basis of the symmetric polynomials. In each degree,
\[
 \langle J_{\lambda_1}\cdots J_{\lambda_d},m_\mu\rangle
 =\begin{cases}\prod_i\lambda_i!,&\lambda=\mu,\\0,&\lambda\ne\mu.
   \end{cases}
\]
The factorials are units in $R_N$, proving perfection in that degree
and the restricted-dual identification. Precomposing the pairing by
$D$ raises one derivative order at a time, with a minus sign.
The derivation rule for $T$ gives the asserted identity.
\end{proof}

Write $\iota_d:A_d^\circ\xrightarrow{\sim}W_{N,d}$ for this
identification. The generating series
\begin{equation}\label{cyjets:eq:generating-state}
 G(u)=\sum_{m\geq0}\frac{u^m}{m!}J_m=e^{uT}J_0
\end{equation}
records all finite jets. It is a formal series of states, not an
additional algebraic state. Taylor's formula gives the useful identity
\begin{equation}\label{cyjets:eq:taylor-pairing}
 \left\langle\prod_{i=1}^dG(u_i),F\right\rangle=F(u_1,\ldots,u_d).
\end{equation}
Every formula below involving these series defines an operation on
algebraic states by taking finitely many auxiliary derivatives at zero.

\section{The full Euler factor}
\label{cyjets:sec:euler}

Consider the three-loop representation space
$\operatorname{End}(\C^d)^3$, with zero potential and arrow weights
$\eps a,\eps b,-\eps h$. Its equivariant cohomology, after this
base change, is $A_d$. This zero-potential coefficient system differs
from the vanishing-cycle system for $\Tr Z[X,Y]$.

Set $\ell=\eps h$ and define
\begin{equation}\label{cyjets:eq:kernels}
 \begin{aligned}
 K_0(s)&=\frac{(s+\eps a)(s+\eps b)(s-\ell)}s,
 &E_0(s)&=K_0(s)^{-1},\\
 E_{\mathrm{pl}}(s)&=\frac{s(s+\ell)}{(s+\eps a)(s+\eps b)},
 &P(s)&=s^2-\ell^2,\\
 r(s)&=\frac{E_{\mathrm{pl}}(s)}{E_{\mathrm{pl}}(-s)}.
 \end{aligned}
\end{equation}
Fractions with nilpotent parameters are expanded in $\eps$ after
inverting the displayed spectral variable. The identity
\begin{equation}\label{cyjets:eq:even-factor}
 E_0(s)=\frac{E_{\mathrm{pl}}(s)}{P(s)},\qquad
 \frac{E_0(s)}{E_0(-s)}=r(s)
\end{equation}
separates the contraction from its exchange ratio.

For $F\in A_d$ and $H\in A_e$, the Hall product is
\begin{equation}\label{cyjets:eq:hall-product}
 (F*H)(X)=
 \sum_{\substack{I\sqcup J=[d+e]\\|I|=d}}
 F(X_I)H(X_J)\prod_{i\in I,\,j\in J}K_0(x_j-x_i).
\end{equation}
There is no factorial divisor in this sum: its summands are the
coordinate flags of the proper invariant-subspace correspondence.
Its spectral cooperation is
\begin{equation}\label{cyjets:eq:hall-split}
 \Delta^0_{d,e}(s)F
 =\prod_{i=1}^d\prod_{j=1}^eE_0(s+y_j-x_i)F(x-s,y),
 \qquad F\in A_{d+e}.
\end{equation}
Here $x-s$ means $(x_1-s,\ldots,x_d-s)$, and expansion at infinity
takes values in $(A_d\otimes_{R_N}A_e)((s^{-1}))$.

\begin{proposition}\label{cyjets:prop:hall}
The Hall product is polynomial and associative. The cooperation is
rationally coassociative: for blocks $x,y,w$ the two iterates use
parameters $s,t,s+t$ and evaluate the input at
$(x-s-t,y-t,w)$. In charge split $(1,1)$,
\begin{equation}\label{cyjets:eq:charge-two}
 1*1=2\bigl((y-x)^2+\eps^2(ab-h^2)\bigr),\qquad
 \Delta^0_{1,1}(s)1=E_0(s+y-x).
\end{equation}
\end{proposition}
\begin{proof}
At a coordinate flag, the three moving arrow spaces contribute
$y-x+\eps a$, $y-x+\eps b$ and $y-x-\ell$; the flag tangent
space contributes $y-x$ in the denominator. Their Euler ratio is
$K_0(y-x)$. The proper pushforward exists before localization.
Algebraically, each possible root-diagonal pole in the shuffle sum
is simple, and its residues cancel between the flags interchanged
across that diagonal. The result is therefore polynomial.

Either bracketing of a triple product sums over ordered three-block
partitions, with one $K_0$ for every ordered cross-block pair. For
the cooperation, both iterates have the common expression
\[
 \prod_{i,j}E_0(s+y_j-x_i)
 \prod_{i,k}E_0(s+t+w_k-x_i)
 \prod_{j,k}E_0(t+w_k-y_j)
 F(x-s-t,y-t,w).
\]
These equalities precede the choice of Laurent expansion. Finally,
$K_0(v)+K_0(-v)$ gives the first formula in
\eqref{cyjets:eq:charge-two}.
\end{proof}

Empty blocks give
\[
 \Delta^0_{0,d}(s)F=F,\qquad
 \Delta^0_{d,0}(s)F=F(x-s).
\]
Thus one empty-block identity includes translation. The cooperation
has not been declared an ordinary Hopf coproduct with two unshifted
counit identities. Nor is the product $*$ the ordinary multiplication
on $W_N$: the latter merely forms polynomial states.

There is an obstruction to replacing these identities by ordinary
counital multiplicativity on the same Hall algebra. Put
\[
 \alpha=\eps^2(ab-h^2),\qquad \gamma=-\eps^3abh,
 \qquad e_m=x^m\in A_1.
\]
Thus $K_0(v)=v^2+\alpha+\gamma/v$ and $e_0$ is the constant
in charge one, distinct from the unit in charge zero.

\begin{proposition}\label{cyjets:prop:no-ordinary-coproduct}
For $N\geq4$, the charge-graded Hall algebra $\bigoplus_dA_d$
admits no $R_N$-linear, charge-preserving, unital algebra map to
its ordinary tensor square whose two-sided counit is projection
onto charge zero. This conclusion does not require coassociativity.
\end{proposition}
\begin{proof}
The shuffle formula gives
\[
 [f,g]_* =2\gamma\frac{f(x)g(y)-g(x)f(y)}{y-x}.
\]
Consequently
\begin{equation}\label{cyjets:eq:primitive-obstruction}
 [e_0,e_3]_* -3[e_1,e_2]_*+\alpha[e_0,e_1]_*
       =\gamma(e_0*e_0).
\end{equation}
Indeed the left side is
$2\gamma((y-x)^2+\alpha)$, and
$e_0*e_0=2((y-x)^2+\alpha)$.
Charge preservation and the two counit identities force every
$e_m$ to be primitive. Commutators of even primitive elements
are primitive in the ordinary tensor product. Applying a
multiplicative coproduct to \eqref{cyjets:eq:primitive-obstruction}
and taking its charge-$(1,1)$ component therefore gives
$0=2\gamma e_0\otimes e_0$. This is impossible in the free
module $A_1\otimes_{R_N}A_1=R_N[x,y]$, since $\gamma\ne0$
for $N\geq4$.
\end{proof}

At $N\leq3$ this particular obstruction vanishes. The proposition
concerns ordinary tensor multiplication and two unshifted counits;
it leaves the spectral cooperation above, with its translated
empty-block identity, intact.

Finite precision will be essential to the lower-truncated field target.
Before nilpotent base change, specialize the arrow weights to $1,1,-2$.
The exchange is
\[
 r(s)=\frac{(s+2)(s-1)^2}{(s-2)(s+1)^2},\qquad
 \operatorname{Res}_{s=2}r(s)\,ds=\frac49.
\]
A finite Laurent polynomial in $s$ has no pole at $s=2$. At precision
$N$, the shifted denominators instead expand into finitely many
powers of $s^{-1}$. Dualizing then changes polynomial differentiation
into the nonnilpotent translation of
Proposition~\ref{cyjets:prop:pairing}.

\section{Fields on polynomial jets}
\label{cyjets:sec:fields}

For lists $u=(u_1,\ldots,u_d)$ and $v=(v_1,\ldots,v_e)$ write
$A_u=\prod_iG(u_i)$ and $B_v=\prod_jG(v_j)$. Define the field
on their generating series by
\begin{equation}\label{cyjets:eq:field}
 Y(A_u,z)B_v=
 \prod_{i,j}E_0(-z-u_i+v_j)
 \prod_iG(z+u_i)\prod_jG(v_j).
\end{equation}
Expand first in the auxiliary variables and in the nilpotent parameter.
Finite auxiliary differentiation defines $Y(a,z)b$ for every
$a,b\in W_N$. Define the exchange operator by the same prescription:
\begin{equation}\label{cyjets:eq:exchange}
 \mathcal S(z)(A_u\otimes B_v)=
 R_{u,v}(z)A_u\otimes B_v,\qquad
 R_{u,v}(z)=\prod_{i,j}r(z+u_i-v_j).
\end{equation}
The notation means coefficient extraction from a generating identity;
on individual jet monomials the operator generally has several terms.

\begin{theorem}\label{cyjets:thm:fields}
The fields have local vacuum $1\in W_{N,0}$ and translation $T$:
\[
 Y(1,z)=\id,\qquad Y(a,z)1=e^{zT}a,\qquad
 Y(Ta,z)=\partial_zY(a,z),\qquad [T,Y(a,z)]=\partial_zY(a,z).
\]
For inputs of charges $d,e$ and total jet order $M$, their values lie
in $W_{N,d+e}((z))$, with pole order at most $2de+N-1+M$.
After pairing with a Hall polynomial, triple products and iterates
are expansions of one rational function. The exchange is unitary and satisfies the spectral
Yang--Baxter equation and the block-fusion identities.
\end{theorem}
\begin{proof}
The empty products give the vacuum identities. Differentiation in
the first auxiliary list agrees with differentiation in $z$.
Differentiation in both lists annihilates every contraction, since
it depends on their differences. These two observations prove the
translation identities after coefficient extraction.

Expand $E_0(q)=\sum_{j=0}^{N-1}c_j\eps^j q^{-2-j}$.
In a product of $de$ contractions, the baseline pole order is $2de$.
A nonzero coefficient has total $\eps$-degree at most $N-1$;
each auxiliary derivative raises the possible pole order by at most
one. The remaining state factor is a power series in $z$.
At each fixed power of $z$, finitely many of its terms contribute.
This proves both the Laurent target and the asserted bound.

For a third list $w$, put $C_w=\prod_kG(w_k)$. The product
$Y(A_u,z_1)Y(B_v,z_2)C_w$ is the indicated expansion of
\begin{align*}
 &\prod_{i,j}E_0(-z_1+z_2-u_i+v_j)
  \prod_{i,k}E_0(-z_1-u_i+w_k)\\
 &\hspace{15mm}\cdot\prod_{j,k}E_0(-z_2-v_j+w_k)
  \prod_iG(z_1+u_i)\prod_jG(z_2+v_j)\prod_kG(w_k).
\end{align*}
The iterate $Y(Y(A_u,z_1-z_2)B_v,z_2)C_w$ gives the same
expression. At fixed auxiliary orders and precision, a finite product
of powers of $z_1$, $z_2$ and $z_1-z_2$ clears the denominators.
This proves rational weak associativity in the respective expansion
regions.

Since $r(q)r(-q)=1$, the exchange is unitary. On three generating
inputs the pairwise exchange factors commute, proving the
Yang--Baxter equation. Compatibility with fusion uses the stronger
identities
\[
 R_{u,v\sqcup w}(z)=R_{u,v}(z)R_{u,w}(z),\qquad
 R_{u\sqcup v,w}(z)=R_{u,w}(z)R_{v,w}(z),
\]
together with substitution of the translated auxiliary coordinates
in a fused block. They follow by partitioning the cross pairs.
Thus fusion compatibility follows directly from the block formula,
in addition to the Yang--Baxter identity.
\end{proof}

Our exchange convention can be checked without terminology. If
$Y_z(a\otimes b)=Y(a,z)b$, then
\begin{equation}\label{cyjets:eq:skew}
 Y_z\mathcal S(z)(A_u\otimes B_v)
       =e^{zT}Y(B_v,-z)A_u.
\end{equation}
Indeed the contraction on the left is
$R_{u,v}(z)\prod E_0(-z-u_i+v_j)=\prod E_0(z+u_i-v_j)$,
and the translated states on both sides agree.

Modulo $\eps^4$, the first coefficients are
\begin{equation}\label{cyjets:eq:cubic}
 \begin{aligned}
 E_0(-z)&=z^{-2}+(h^2-ab)\eps^2z^{-4}-abh\eps^3z^{-5},\\
 r(z)&=1+2abh\eps^3z^{-3}.
 \end{aligned}
\end{equation}
They follow by inverting
$K_0(s)=s^2+\eps^2(ab-h^2)-\eps^3abh/s$.
The even factor $P^{-1}$ changes the leading contraction from $1$
to $z^{-2}$ while leaving the exchange ratio fixed. Exchange data
alone therefore do not determine the field.

\begin{theorem}[The restricted transpose]\label{cyjets:thm:transpose}
For finite-jet functionals $\alpha\in A_d^\circ$ and
$\beta\in A_e^\circ$,
\begin{equation}\label{cyjets:eq:transpose}
 Y(\iota_d\alpha,z)\iota_e\beta
 =\iota_{d+e}\bigl((\Delta^0_{d,e}(-z))^\circ
                                    (\alpha\otimes\beta)\bigr).
\end{equation}
The right side is defined coefficientwise on polynomial tests and
belongs to $W_{N,d+e}((z))$. Its exchange is the transpose of the
inverse Hall multiplier at the reflected coordinate $-z$.
\end{theorem}
\begin{proof}
Pair \eqref{cyjets:eq:field} with $F\in A_{d+e}$. Equation
\eqref{cyjets:eq:taylor-pairing} gives
\[
 \prod_{i,j}E_0(-z-u_i+v_j)F(z+u_1,\ldots,z+u_d,v_1,\ldots,v_e).
\]
This is precisely \eqref{cyjets:eq:hall-split} at $s=-z$,
evaluated on $x=u,y=v$. Finite auxiliary differentiation proves
the equality. A fixed Laurent coefficient uses finitely many
derivatives of $F$, so it is in the restricted dual; the theorem's
lower bound is supplied by Theorem~\ref{cyjets:thm:fields}.
Finally the Hall multiplier is $\prod r(s+y_j-x_i)$, whose inverse
at $s=-z$ is $\prod r(z+x_i-y_j)$ by unitarity. This is
\eqref{cyjets:eq:exchange} under the derivative pairing.
\end{proof}

Give roots and $\eps$ degree two and $a,b$ degree zero. A root-degree
$m$ Hall polynomial in charge $d$ has shifted degree $2m-2d^2$.
Its restricted dual state has degree $2d^2-2m$. With $\deg z=2$
and $\deg T=-2$, the field comparison is homogeneous. Ordinary
polynomial multiplication on $W_N$ does not have degree zero for
this charge-dependent shift. This grading is separate from the
cochain grading, which remains zero.

The completions also have distinct meanings. For charge-completed
states the field target is
\[
 \prod_{d\geq0}W_{N,d}((z));
\]
at a fixed total charge only finitely many input charge splittings
contribute. No uniform pole bound over all charges is imposed.
For a formal parameter retain the compatible system
$\varprojlim_NW_N((z))$, with a bound at each precision.
The cooperation on polynomials has not been extended to every
root-adic series: translation by the nonnilpotent variable $s$
does not preserve the root-adic topology. Its coefficientwise
restricted transpose is the map proved above.

\section{Moving the centre of evaluation}
\label{cyjets:sec:moving}

Extend coefficients to $R=R_N[a_{\mathrm{fr}},\beta]$ and put
$\lambda=a_{\mathrm{fr}}+\eps b$. Here $a_{\mathrm{fr}}$ is the
moving centre parameter; $\beta$ scales its derivative. The connection
on charge-$d$ polynomials is
\begin{equation}\label{cyjets:eq:connection}
 \nabla_d=\beta\left(\partial_{a_{\mathrm{fr}}}
                              +\sum_i\partial_{x_i}\right).
\end{equation}
It obeys the Leibniz rule over the coefficient derivation
$\beta\partial_{a_{\mathrm{fr}}}$. On a tensor product over $R$,
that coefficient derivative occurs once, while every root derivative
occurs once.

\begin{proposition}\label{cyjets:prop:moving}
The Hall product and cooperation are horizontal for this connection.
The finite derivative evaluations at $x_i=\lambda$, indexed by
nondecreasing derivative orders, form a horizontal basis of the
continuous dual of the completion along that centre.
Their identification with $W_{N,d}\otimes_{R_N}R$ makes the field
transpose horizontal.
\end{proposition}
\begin{proof}
Every Euler factor depends on root differences, so the sum of its
root derivatives is zero. Translation of one root block by $s$
commutes with simultaneous root differentiation. This proves the
first assertion, including the flag sum in the product.

For $\xi_i=x_i-\lambda$, one has $\nabla_d\xi_i=0$.
The completion in the $\xi_i$ is therefore preserved. If
\[
 v_{\mathbf m,\lambda}(F)=
 \left.\partial_{x_1}^{m_1}\cdots\partial_{x_d}^{m_d}F
                                                    \right|_{x_i=\lambda},
\]
the chain rule gives
\[
 \beta\partial_{a_{\mathrm{fr}}}
       \langle v_{\mathbf m,\lambda},F\rangle
       =\langle v_{\mathbf m,\lambda},\nabla_dF\rangle.
\]
This is horizontality for the dual connection. The assignment
$v_{\mathbf m,\lambda}\mapsto\prod_iJ_{m_i}$ reduces the transpose
proof to evaluation of $F$ at $(\lambda+z+u,\lambda+v)$.
The contraction factors are unchanged.
\end{proof}

For the test $F(x)=x$, both sides of the displayed chain rule equal
$\beta$. This test detects omission of the coefficient derivative.
The moving-centre identification uses new evaluation functionals.
When $\lambda$ is nonnilpotent, the expression $e^{\lambda T}J_0$
is an infinite series and does not give an algebraic vector in $W_N$.
It cannot be substituted for that change of coefficient completion.

\section{Why a singular current needs an infinite orbit}
\label{cyjets:sec:orbits}

The jet states above have independent translates. A pole forces this
phenomenon in every translation-equivariant field system over a field
of characteristic zero, regardless of its Hall origin.

\begin{lemma}\label{cyjets:lem:orbit}
Let $V$ be a $k$-vector space with fields $Y(a,z)b\in V((z))$ satisfying
$Y(Ta,z)b=\partial_zY(a,z)b$. If $Y(a,z)b$ has a nonzero
negative-power term, the states $a,Ta,T^2a,\ldots$ are linearly
independent.
\end{lemma}
\begin{proof}
Write the lowest term of $Y(a,z)b$ as $v z^{-p}$ with $p>0$
and $v\ne0$. Suppose $\sum_{j=0}^m c_jT^ja=0$ with $c_m\ne0$.
Applying its field to $b$ gives a constant-coefficient differential
equation. Its lowest power is $z^{-p-m}$, with coefficient
\[
 c_m(-p)(-p-1)\cdots(-p-m+1)v\ne0.
\]
The other derivatives have larger lowest powers, a contradiction.
\end{proof}

Over $R_N$, the same conclusion holds for a flat state module when
the pole remains nonzero modulo $\eps$. For a finite relation, take
the least $\eps$-adic order of its coefficients. Flatness identifies
the corresponding associated-graded piece with the reduction modulo
$\eps$, where the lemma applies. The hypothesis on the reduced pole
is necessary: over $k[\eps]/(\eps^2)$ the state $\eps a$ is
annihilated by $\eps$, even if its field has a nonzero polar term.

Lower truncation must also be interpreted separately for each second
input. For the level-one Heisenberg module put
\[
 V=k[q_1,q_2,\ldots],\qquad
 a_{-n}=q_n,\quad a_n=n\partial_{q_n}\ (n>0),\quad a_0=0,
 \qquad J=q_1.
\]
The field is $Y(J,z)=\sum_{n\in\Z}a_nz^{-n-1}$ and the vacuum
is $1$. Translation on states is the derivation
$Tq_n=nq_{n+1}$. In particular,
\begin{equation}\label{cyjets:eq:unbounded-poles}
 T^rJ=r!q_{r+1},\qquad
 a_{r+1}T^rJ=(r+1)!1.
\end{equation}
The fixed field $Y(J,z)$ therefore has no lower Laurent bound
uniform over its inputs. Its correct target is
$\Hom_k(V,V((z)))$; writing it as an element of the literal space
$\End_k(V)((z))$ would impose a uniform bound. Every polynomial
input still has a finite lower bound.

Scalar completion introduces a second failure of uniformity, now
for one fixed pair of completed inputs. Extend this Heisenberg module
to $V[[\eps]]$ and consider the state
\[
 a(\eps)=\sum_{n\geq0}\frac{(-\eps)^n}{(n+1)!}T^nJ.
\]
Its vacuum-valued contraction with $J$ is
\begin{equation}\label{cyjets:eq:completed-poles}
 [1]Y(a(\eps),z)J=\sum_{n\geq0}\eps^n z^{-n-2},
\end{equation}
where $[1]$ extracts the constant polynomial state. To obtain this
formula, differentiate the contraction $z^{-2}$ exactly $n$ times:
$\partial_z^nz^{-2}=(-1)^n(n+1)!z^{-n-2}$.
The resulting series has no finite lower bound in $z$, although
its reduction modulo every $\eps^r$ does.

For an $\eps$-adically complete module $M$, the appropriate recipient is
\begin{equation}\label{cyjets:eq:restricted-laurent}
 \mathcal L_\eps(M)=\varprojlim_r(M/\eps^rM)((z)).
\end{equation}
Equivalently, it consists of series $\sum_{j\in\Z}m_jz^j$ such
that for each $r$ all sufficiently negative $j$ satisfy
$m_j\in\eps^rM$. A compatible family of finite-precision Laurent
series determines its coefficients in $M$ and has exactly this
property, proving the equivalence. Equation
\eqref{cyjets:eq:completed-poles} belongs to this recipient, whereas
it does not belong to $V[[\eps]]((z))$. Completion of the states
therefore requires completion of the Laurent recipient as well.

\subsection{A quotient and its dual have opposite directions}

A finite Hamiltonian quotient gives another useful test of what a
jet truncation preserves. Let $\mathfrak m=(x,y)\subset k[[x,y]]$
and give $\mathfrak g=\mathfrak m^2$ the Poisson bracket
$\{x,y\}=1$. For $d\geq2$, the subspace
$\mathfrak m^{d+1}$ is a Lie ideal in $\mathfrak g$; hence
$\mathfrak g_d=\mathfrak m^2/\mathfrak m^{d+1}$ is a finite
Lie algebra. Let $\rho_{ab}$ extract the coefficient of $x^ay^b$,
where $a,b\geq0$ and $a+b\geq2$. These span the continuous dual
$\mathfrak g^\circ$.

The coadjoint action is defined by
$(f\cdot\rho)(g)=-\rho(\{f,g\})$. It gives
\begin{equation}\label{cyjets:eq:coadjoint}
 x^iy^j\cdot\rho_{ab}
 =\bigl(j(a+1)-i(b+1)\bigr)\rho_{a-i+1,b-j+1},
\end{equation}
with the result zero when the indicated functional is outside
$\mathfrak g^\circ$. To verify the coefficient, pair with $x^py^q$:
the only possible indices are $p=a-i+1$, $q=b-j+1$, and the
coadjoint sign changes $iq-jp$ to the displayed number.

Each shifted cotangent Lie algebra
$\mathfrak g_d\ltimes\mathfrak g_d^\vee[-1]$ is well defined.
However, simultaneous projection of both factors by total degree
need not give its transition maps. Already
\begin{equation}\label{cyjets:eq:cotangent-witness}
 x^3\cdot\rho_{21}=-6\rho_{02}.
\end{equation}
Projection to degree at most two kills both inputs and preserves
the nonzero output. It is therefore not a Lie map on the cotangent
algebra. The quotient $\mathfrak g\to\mathfrak g_2$ instead
dualizes to an inclusion $\mathfrak g_2^\vee\to\mathfrak g^\circ$.
A compatible family of cotangent theories must supply maps with
their actual variance; it cannot use the simultaneous projection.

\section{Perfect differential modules and their pro-duals}
\label{cyjets:sec:perfect}

Let $X=\mathbb A^1_k$ with coordinate $u$ and differential-operator
algebra
\[
 D=k\langle u,\partial\rangle/(\partial u-u\partial-1).
\]
A state space $(V,T)$ defines a right differential module
$\omega_X\otimes_kV$ by multiplication by functions and
\begin{equation}\label{cyjets:eq:right-current}
 (f\,du\otimes v)\partial
       =-f'\,du\otimes v-f\,du\otimes Tv.
\end{equation}
This formula satisfies the Weyl relation. If $k[T]a$ is a free
translation orbit, the right module $\omega_X\otimes k[T]a$
is free of rank one over $D$: the generator $du\otimes a$
sends $\partial^m$ to $(-1)^mdu\otimes T^ma$. The formula
with polynomial coefficients is triangular in the derivative order,
so it gives an isomorphism. Its characteristic variety is the full
two-dimensional cotangent bundle. In particular it is not holonomic,
which on a curve would require characteristic dimension one.

The finite object appropriate to such a coefficient is therefore
a finite complex of free differential modules, rather than a finite
span of current modes. Let $X$ now be a smooth affine coordinate
chart of dimension $n$, and let $D_X=\Gamma(X,\mathcal D_X)$.
A perfect right $D_X$-module is a retract of a finite complex of
finite free right modules. Denote their stable category by
$\Perf_r(D_X)$. All module categories in this section are derived
categories with their mapping spaces, so homotopies are retained.

A coordinate volume form converts left differential modules to right
ones. The formal transpose is the anti-involution
\[
 t(f)=f,\qquad t(\partial_i)=-\partial_i,\qquad
 t(PQ)=t(Q)t(P).
\]
For $P\in\Perf_r(D_X)$ define
\begin{equation}\label{cyjets:eq:dagger}
 P^\dagger=\bigl(\RHom_{D_X}(P,D_X)\bigr)^t[n].
\end{equation}
The superscript converts the left module into a right module; the
shift is the dimension normalization. Intrinsically the canonical
line supplies the left--right conversion.

\begin{lemma}\label{cyjets:lem:perfect}
The functor $(-)^\dagger$ is an involutive anti-equivalence on
$\Perf_r(D_X)$. It commutes with open restriction and with external
products on products of smooth charts.
\end{lemma}
\begin{proof}
Evaluation is an isomorphism from a finite projective module to
its double $D_X$-dual. Finite complexes and retracts preserve it.
Duality reverses a cochain shift, so the two shifts by $n$ cancel
in the double dual. Localization gives the open-restriction identity.
For an external product, evaluation gives the comparison on free
modules over the two operator algebras, and finite cones and retracts
extend it to perfect modules. The dimensions add. Canonical-line
conversion makes these comparisons agree on intersections of charts.
\end{proof}

For an associative ring, compact objects in its derived module
category are exactly the perfect modules, and every module is a
filtered colimit of perfect modules
\cite{LurieHA}*{Proposition 7.2.4.2}. The perfect left--right
duality is also expressed in \cite{LurieHA}*{Proposition 7.2.4.7}.
No holonomicity condition occurs in these statements. Thus write
\[
 M\simeq\varinjlim_\alpha P_\alpha.
\]
The maps in this presentation need not be inclusions of submodules.
Dualize the presentation as a diagram:
\begin{equation}\label{cyjets:eq:sharp}
 M^\sharp=\{P_\alpha^\dagger\}_\alpha
             \in\operatorname{Pro}\Perf_r(D_X).
\end{equation}
For a pro-perfect diagram $Q=\{Q_\beta\}$ set
$\mathcal UQ=\varinjlim_\beta Q_\beta^\dagger$.

\begin{theorem}[Full coefficient duality]\label{cyjets:thm:pro}
These constructions give inverse anti-equivalences
\begin{equation}\label{cyjets:eq:pro-equivalence}
 \Mod_r(D_X)^{\op}\simeq\operatorname{Pro}\Perf_r(D_X).
\end{equation}
In particular $M^\sharp$ is independent of the chosen perfect
presentation.
\end{theorem}
\begin{proof}
Let $M=\varinjlim_\alpha P_\alpha$ and
$N=\varinjlim_\beta Q_\beta$. The pro-category mapping formula,
perfect biduality and compactness give
\begin{align*}
 \Map_{\operatorname{Pro}}(N^\sharp,M^\sharp)
 &\simeq\varprojlim_\alpha\varinjlim_\beta
                   \Map(Q_\beta^\dagger,P_\alpha^\dagger)\\
 &\simeq\varprojlim_\alpha\varinjlim_\beta
                   \Map(P_\alpha,Q_\beta)\\
 &\simeq\Map(M,N).
\end{align*}
The equivalence respects composition. It proves full faithfulness
and makes the presentation immaterial. Conversely, the same formula
identifies any $\{Q_\beta\}$ with
$(\varinjlim_\beta Q_\beta^\dagger)^\sharp$. Evaluation supplies
both inverse comparisons and their naturality.
\end{proof}

\subsection{The diagram cannot be replaced by its limit}

On the affine line, the localized free right module has the perfect
presentation
\begin{equation}\label{cyjets:eq:localization}
 D[u^{-1}]=\varinjlim
       \bigl(D\xrightarrow{L_u}D\xrightarrow{L_u}D\longrightarrow\cdots\bigr),
\end{equation}
where $L_u$ is left multiplication by $u$. The stage-$m$ map
sends $a$ to $u^{-m}a$. Normal form in powers of $\partial$
shows that these images exhaust the localization. Its dual is
\begin{equation}\label{cyjets:eq:dual-tower}
 D[u^{-1}]^\sharp=
       \{\cdots\xrightarrow{L_u}D[1]\xrightarrow{L_u}D[1]\}.
\end{equation}
Indeed precomposition by $L_u$ sends a left-module value $b$ to
$bu$, and transpose converts this to $L_u$ on a right module.
Every transition to a fixed stage is a nonzero power of $L_u$.
The pro-object is therefore nonzero.

Its underived inverse limit is zero. The first component of a
compatible sequence belongs to $\bigcap_m u^mD=0$, since each
coefficient in its finite normal form would be divisible by every
power of $u$. The same argument applies at every stage. This
calculation does not assert vanishing of a derived inverse limit.

The vacuum coefficient has a different dual. The right-module
resolution
\[
 0\longrightarrow D\xrightarrow{L_\partial}D
                  \longrightarrow\omega_X\longrightarrow0
\]
and \eqref{cyjets:eq:dagger} give
$\omega_X^\dagger\simeq\omega_X$, whereas $D^\dagger=D[1]$.
The same category therefore contains the vacuum and the free
current, with different resolutions and shifts.

\subsection{Finite perfect pieces of the polynomial current}

Return to $P=k[J_0,J_1,\ldots]$ with $TJ_m=J_{m+1}$, and give
$J_m$ weight $m+1$. Each weight space is finite-dimensional.
The derivation $T$ is injective on the positive-weight part $P_+$.
To see this, take a nonconstant polynomial whose highest occurring
variable is $J_M$. The coefficient of $J_{M+1}$ in its derivative
is $\partial F/\partial J_M\ne0$.

Choose a complement $Q_w$ to $T(P_+\cap P_{w-1})$ in $P_w$,
and put $Q=\bigoplus_{w>0}Q_w$. Induction on weight, using
injectivity, gives
\begin{equation}\label{cyjets:eq:free-current-presentation}
 P=k1\oplus k[T]\otimes_kQ.
\end{equation}
For example the Hilbert series of $Q$ is
\[
 \sum_{w>0}\dim(Q_w)t^w
   =(1-t)\left(\prod_{n\geq1}(1-t^n)^{-1}-1\right).
\]
Consequently the current module of $P$ is a filtered union of
$\omega_X$ and finite sums of free copies of $D$, one for each
chosen finite set of homogeneous basis elements of $Q$.
Every such piece retains the entire translation orbit of each
generator. The finite presentation is by differential modules,
not by a cutoff on the energy or derivative order.

\section{Derived normal support}
\label{cyjets:sec:normal-support}

The pro-dual retains the full coefficient. A collision also singles out
the coefficients supported where a normal coordinate vanishes. This
support condition must be checked after derived scalar reduction:
ordinary reduction can omit torsion cohomology. We examine it on one
normal line, where both reconstruction and convolution are explicit.

In this section put $R=k[[\eps]]$, $R_r=R/(\eps^r)$, and
\[
 D_R=R\langle u,\partial\rangle/(\partial u-u\partial-1).
\]
The formal scalar $\eps$ has cochain degree zero and is independent
of the normal coordinate $u$.
For a complex $M$ with an $R$-action write
\[
 M_r=M\otimes_R^{\mathbb L}R_r,\qquad
 \Lambda M=R\!\varprojlim_r M_r.
\]
Call $M$ complete when $M\to\Lambda M$ is an equivalence.
Restriction $M[u^{-1}]$ uses flat localization of the underlying
$R[u]$-module and its induced differential action. It does not treat
$u$ as central in the Weyl algebra.

Let $\mathcal G_\eps$ be the derived category of complete scalar
$R$-complexes. Let $\mathcal C_\eps$ be the full category of complete
right $D_R$-complexes satisfying
\begin{equation}\label{cyjets:eq:derived-support}
 M_r[u^{-1}]\simeq0\qquad\text{for every }r\geq1.
\end{equation}
This is the derived adic normal support condition. It imposes no
boundedness or scalar-flatness assumption. The condition at $r=1$
already suffices: it says that multiplication by $\eps$ is an
equivalence on $M[u^{-1}]$, and hence that all its derived reductions
vanish. For an $R$-flat module the condition becomes
\[
 \text{for every }m\in M\text{ and }r\geq1\text{ there is }n
       \text{ with }mu^n\in\eps^rM.
\]
Without flatness, this elementwise condition controls only the
ordinary quotients and is not the definition in
\eqref{cyjets:eq:derived-support}.

The basic supported right module is the principal-parts module
\begin{equation}\label{cyjets:eq:normal-parts}
 P=R[u,u^{-1}]/R[u]=\bigoplus_{m\geq0}R\beta_m,
 \qquad \beta_m=[u^{-m-1}].
\end{equation}
Right differentiation is minus ordinary differentiation. Thus
\[
 \beta_0u=0,\qquad \beta_mu=\beta_{m-1}\ (m>0),\qquad
 \beta_m\partial=(m+1)\beta_{m+1}.
\]
The map $1\mapsto\beta_0$ identifies $P=D_R/uD_R$ and gives
the resolution
\begin{equation}\label{cyjets:eq:normal-resolution}
 0\longrightarrow D_R\xrightarrow{L_u}D_R
                       \longrightarrow P\longrightarrow0.
\end{equation}
Ordered Weyl monomials prove injectivity of $L_u$.

\begin{theorem}[Normal reconstruction]\label{cyjets:thm:normal}
The functors
\[
 I(E)=\Lambda(P\otimes_R^{\mathbb L}E),\qquad
 K(M)=\RHom_{D_R}(P,M)
\]
give inverse equivalences
$I:\mathcal G_\eps\simeq\mathcal C_\eps:K$.
For a right differential module complex,
\begin{equation}\label{cyjets:eq:normal-ground}
 K(M)=[M\xrightarrow{\,u\,}M]_{[0,1]}.
\end{equation}
The last display is a complex of scalar modules; right multiplication
by $u$ need not be a differential-module map.
\end{theorem}
\begin{proof}
First work over any commutative $k$-algebra $B$, with the analogous
Weyl algebra $D_B$ and $P_B=B[u,u^{-1}]/B[u]$. If a right
$D_B$-module $V$ is supported at $u=0$, multiplication by $u$
is locally nilpotent. For $v\in V$ the finite sum
\[
 \pi(v)=\sum_{j\geq0}\frac{(-1)^j}{j!}(vu^j)\partial^j
\]
belongs to $\ker u$. Indeed
$\partial^ju=u\partial^j+j\partial^{j-1}$ makes consecutive
terms of $\pi(v)u$ cancel. The same binomial cancellation gives
\[
 v=\sum_{m\geq0}\frac{\pi(vu^m)\partial^m}{m!}.
\]
Consequently
\begin{equation}\label{cyjets:eq:normal-evaluation}
 P_B\otimes_B\ker u\xrightarrow{\sim}V,
 \qquad \beta_m\otimes e\longmapsto e\partial^m/m!.
\end{equation}
Surjectivity follows from the preceding identity. In a finite
relation, applying the largest occurring power of $u$ extracts
the highest coefficient, proving injectivity. This also proves
that $u:V\to V$ is surjective.

Resolution \eqref{cyjets:eq:normal-resolution} gives
\eqref{cyjets:eq:normal-ground}. For a complex with supported
cohomology, its long exact sequence therefore identifies
$H^qK(V)$ with $\ker(u:H^qV\to H^qV)$. Since $P_B$ is
$B$-free, \eqref{cyjets:eq:normal-evaluation} proves that the
adjunction counit $P_B\otimes_B^{\mathbb L}K(V)\to V$ is
an equivalence in every degree. The unit is an equivalence because
$u:P_B\to P_B$ is a split surjection with kernel $B\beta_0$.
No truncation of the complex was used.

Derived completion preserves every derived reduction:
$(\Lambda X)_r\simeq X_r$. To check this, tensor the defining
limit with the perfect $R$-complex $R_r$. For $n\geq r$ the
resulting tower has a constant copy of $X_r$ and a shifted copy
whose transitions multiply by $\eps$. Any $r$ successive
transitions on the latter copy vanish, so its homotopy limit is
zero. Completion is the reflection onto complete complexes: its
kernel consists of complexes on which $\eps$ is invertible,
and maps from such a complex to a complete one vanish.

The finite free resolution of $P$ shows that $K$ commutes with
reduction and completion. Tensor--Hom adjunction and the reflection
just described therefore give $I\dashv K$ on complete objects.
Its unit and counit reduce to the equivalences over $B=R_r$
proved above. Complete objects are limits of their reductions,
so these maps are equivalences over $R$ as well.
\end{proof}

\subsection{Additive convolution}

There is a geometric tensor preserving this support. For addition
$a:\mathbb A_x^1\times\mathbb A_y^1\to\mathbb A_t^1$,
$t=x+y$, let $a_+$ denote right differential-module direct image,
computed here by the relative right de Rham complex. Define
\[
 M*_+N=\Lambda a_+(M\boxtimes_R^{\mathbb L}N),\qquad
 E\widehat\otimes_R F=\Lambda(E\otimes_R^{\mathbb L}F).
\]

\begin{proposition}\label{cyjets:prop:normal-convolution}
There are natural coherent equivalences
\[
 I(E)*_+I(F)\simeq I(E\widehat\otimes_R F).
\]
They make normal reconstruction symmetric monoidal, with convolution
unit $\widehat P=I(R)$. Relative convolution over an algebra $I(B)$
is given by
\begin{equation}\label{cyjets:eq:normal-relative}
 I(E)*_{I(B)}^{\mathbb L}I(F)
       \simeq I\bigl(\Lambda(E\otimes_B^{\mathbb L}F)\bigr)
\end{equation}
for right and left $B$-modules $E,F$ in complete scalar complexes.
\end{proposition}
\begin{proof}
At precision $r$, change coordinates to $t=x+y$, $s=y$.
Since $\partial_x=\partial_t$ and
$\partial_y=\partial_t+\partial_s$, divided derivatives give
\[
 \beta_i(x)\beta_j(y)\longmapsto
 \sum_{q=0}^j\binom{i+q}{q}\beta_{i+q}(t)\beta_{j-q}(s).
\]
The relative de Rham complex
$[P_s\xrightarrow{\partial_s}P_s]_{[-1,0]}$ has injective
differential and cokernel $R_r\beta_0$. Thus direct image retains
the $q=j$ term:
\begin{equation}\label{cyjets:eq:normal-binomial}
 \beta_i\beta_j\longmapsto\binom{i+j}{i}\beta_{i+j}.
\end{equation}
There is no cohomological shift. This split computation remains
valid after tensoring with any scalar complexes. Associativity
on three modes is
\[
 \binom{i+j}{i}\binom{i+j+k}{i+j}
 =\frac{(i+j+k)!}{i!j!k!}
 =\binom{j+k}{j}\binom{i+j+k}{i}.
\]
The coherent comparisons come from iterated addition and relative
de Rham direct image. Passing through all derived reductions proves
the completed assertion. The two-sided bar for relative convolution
has terms $E\widehat\otimes B^{\widehat\otimes n}
\widehat\otimes F$; its realization gives
\eqref{cyjets:eq:normal-relative}.
\end{proof}

This convolution retains the normal coordinate. Balanced tensor over
the Weyl algebra gives a different output. With formal transpose
$u\mapsto u$, $\partial\mapsto-\partial$, the same resolution gives
\[
 \RHom_{D_R}(P,P)\simeq R,\qquad
 P\otimes_{D_R}^{\mathbb L}P^t\simeq R[1].
\]
The lowering operator is surjective with kernel $R\beta_0$.
Its Hom complex has degrees $0,1$, while its tensor complex has
degrees $-1,0$. These calculations explain both the shift and the
scalar output of the balanced tensor. Identifying additive convolution
with a collision convolution requires a comparison of their actual
correspondences and units.

Derived support is stable under the relative operation even when
scalar flatness is lost. For example, let $B=\widehat{R[t]}$ be
the $\eps$-adic completion, and give $E_0=R$ and $E_\eps=R$
the actions $t=0$ and $t=\eps$. The exact free resolution of
$E_0$ by multiplication by $t$ gives
\[
 E_0\otimes_B^{\mathbb L}E_\eps
       =[R\xrightarrow{\eps}R]_{[-1,0]}\simeq k.
\]
Both inputs are scalar-flat; the output is complete and nonflat.
Its derived reduction has a copy of $k$ in each of degrees $-1,0$.
Applying $I$ retains normal support in both degrees.

Scalar perfection alone is also not stable under arbitrary relative
tensor products. For $B=R[\eta]/(\eta^2)$ and the augmentation
module $R$, the resolution with every negative differential equal
to multiplication by $\eta$ gives
\[
 H^{-n}(R\otimes_B^{\mathbb L}R)=R\qquad(n\geq0).
\]
Both scalar inputs are perfect, but their relative tensor is not.
Perfection of the right module over $B$, together with perfection
of the left module over $R$, is a sufficient stronger hypothesis:
a finite $B$-cell presentation then tensors to a finite $R$-cell
presentation.

\subsection{Limits and ordinary realized duals}

Normal reconstruction describes the supported limits intrinsically:
\[
 \varprojlim_{\mathcal C_\eps}M_i
       \simeq I\left(\varprojlim_iK(M_i)\right).
\]
They need not agree with limits in all complete differential modules.
Indeed
\[
 \widehat P=\left\{\sum_{m\geq0}a_m\beta_m:
       a_m\in R,\ a_m\longrightarrow0\text{ adically}\right\}
\]
is flat, complete and derived-supported. Its ambient countable product
$W=\prod_{i\geq0}\widehat P$ is flat and complete, with
$W_1=\prod_iP_1$. The element
\begin{equation}\label{cyjets:eq:diagonal-product}
 e=(\beta_0,\beta_1,\beta_2,\ldots)\in W_1
\end{equation}
is killed by no power of $u$: the $n$th coordinate of $eu^n$
is $\beta_0\ne0$. Hence $W_1[u^{-1}]\ne0$ and
$W\notin\mathcal C_\eps$.

In matrix coordinates $a_{i,m}$, the supported product
$I(\prod_iR)$ consists of those elements of $W$ satisfying
\[
 \text{for each }r\text{ there is }M_r\text{ such that }
 a_{i,m}\in\eps^rR\quad\text{for every }i\text{ and }m\geq M_r.
\]
The ambient product imposes the bound separately in each row.
The diagonal element \eqref{cyjets:eq:diagonal-product} distinguishes
the two. This is a failure of normal support, even though the tower
of finite products has split surjective transitions.

The same example obstructs an ordinary realized coefficient dual.
Put $\widehat D_R=\Lambda D_R$ and define
\[
 \mathbb D_\eps(M)=\RHom_{D_R}(M,\widehat D_R)^t[1].
\]
This takes values in ordinary complete differential-module complexes.
Dualizing \eqref{cyjets:eq:normal-resolution} yields an injective
right-multiplication map $R_u$ on $\widehat D_R$, now regarded as
a left module. Ordered monomials $\partial^ju^i$ identify its
cokernel with the completed left delta module. Thus
\[
 \RHom_{D_R}(P,\widehat D_R)\simeq\widehat P^{\,t}[-1].
\]
For the flat, complete, supported module
\[
 M=\Lambda\left(\bigoplus_{i\geq0}P\right)
\]
the reflection property of completion and mapping out of a direct
sum now give
\begin{equation}\label{cyjets:eq:ordinary-dual-failure}
 \mathbb D_\eps(M)\simeq\prod_{i\geq0}\widehat P=W
                       \notin\mathcal C_\eps.
\end{equation}
For a perfect scalar ground $E$, finite evaluation instead gives
\[
 \mathbb D_\eps(I(E))\simeq I(\RHom_R(E,R)).
\]
This is a sufficient domain for ordinary supported duality. The
relative-tensor example above shows why scalar perfection alone
does not make that domain stable under every relative fusion.
The perfect-presentation argument of Theorem~\ref{cyjets:thm:pro}
also applies to $D_R$. It retains a diagram. The failure in
\eqref{cyjets:eq:ordinary-dual-failure} concerns the ordinarily realized dual.

\section{Transposing the full bar coalgebra}
\label{cyjets:sec:bar}

Coefficient duality reverses maps. It becomes a comparison of
operations only after the tensor and geometric functors on the
dual coefficients have been specified. On a smooth scheme use
the affine construction above and open descent to define
$\Dcal^\sharp(X)$. The canonical-line identifications agree on
overlaps, and taking opposite categories preserves the descent
limit. Thus the derived category $\Dcal(X)$ of right differential
modules satisfies
$\Dcal(X)^{\op}\simeq\Dcal^\sharp(X)$.

For a given functor $F:\Dcal(X)\to\Dcal(Y)$ define
\begin{equation}\label{cyjets:eq:transport}
 F^\sharp=\sharp_Y\circ F\circ\mathcal U_X.
\end{equation}
In particular, a collision pullback $\delta^!$ gives
$\delta^\dagger=\sharp\delta^!\mathcal U$, and an open direct
image $j_*$ gives $j_{!,\operatorname{pro}}=\sharp j_*\mathcal U$.
The latter takes values in pro-perfect coefficients; it has not
been identified with an ordinary $j_!$ on all differential modules.
The localized module in \eqref{cyjets:eq:dual-tower} explains why
retaining the pro-diagram matters.

Suppose now that $\Ccal$ is a stable category of coherent coefficient
systems on smooth collision spaces $X_I$, indexed by finite label
sets $I$. For an ordered decomposition $I=A\sqcup B$, let
$U_{A|B}\subset X_I$ be the open locus where points in different
blocks are distinct, and let $j_{A|B}$ be its inclusion.
Collisions within one block remain allowed. Assume that external
products and the collision maps define an associative convolution
\[
 (M\star N)_I=
 \bigoplus_{I=A\sqcup B}j_{A|B,*}
           (M_A\boxtimes N_B)|_{U_{A|B}}
\]
with unit $K_\varnothing=k$, $K_I=0$ for $I\ne\varnothing$.
The collision equivalences and their higher compatibilities are
part of a coefficient system. All ordered subset decompositions
occur in this convolution. Keeping only consecutive cuts of an
ordered word would define another tensor.

Componentwise duality gives a category $\Ccal^\sharp$ and an
anti-equivalence $\Ccal^{\op}\simeq\Ccal^\sharp$. Define
\begin{equation}\label{cyjets:eq:dual-tensor}
 E\star_\sharp F=(\mathcal UE\star\mathcal UF)^\sharp.
\end{equation}
Equivalently its component at $I$ is
\[
 \bigoplus_{I=A\sqcup B}j_{A|B,!,\operatorname{pro}}
                 (E_A\boxtimes^\sharp F_B)|_{U_{A|B}},
\]
where $\boxtimes^\sharp$ is the two-index external product of
perfect diagrams. The sum is finite. Associativity, its coherent
homotopies and the unit $K^\sharp=K$ are transported from $\star$.

\begin{proposition}\label{cyjets:prop:operations}
For finite lists of coefficient objects there is a natural equivalence
\[
 \Map_{\Ccal}\left(\mathop{\star}_iM_i,\mathop{\star}_jN_j\right)
 \simeq
 \Map_{\Ccal^\sharp}
       \left(\mathop{\star_\sharp}_jN_j^\sharp,
             \mathop{\star_\sharp}_iM_i^\sharp\right).
\]
It reverses composition and preserves identities and all coherent
homotopies. Collision maps are compared by
\eqref{cyjets:eq:transport}.
\end{proposition}
\begin{proof}
Apply coefficient anti-equivalence to the two tensor objects.
Equation \eqref{cyjets:eq:dual-tensor} identifies their images.
Naturality of the mapping equivalence gives composition and every
collision diagram, including their specified homotopies.
\end{proof}

Let $A$ be an augmented associative algebra for $\star$, with
$A_\varnothing=k$. Its augmentation ideal $\bar A$ has zero
empty coefficient. A normalized bar word of length $r$ on $I$
uses an ordered decomposition into $r$ nonempty blocks. Hence
$r\leq|I|$. This finiteness is in configuration support and makes
no assertion about the dimension of a current coefficient.

Write $B_\star A$ for its full normalized bar coalgebra, with
suspensions of $\bar A$, deconcatenation coproduct and complete
bar differential. For a coherent associative presentation, that
differential includes all higher products. Define
\begin{equation}\label{cyjets:eq:bar-dual}
 \mathcal K_\sharp A=(B_\star A)^\sharp.
\end{equation}
Its multiplication is the explicit transpose
\[
 \mathcal K_\sharp A\star_\sharp\mathcal K_\sharp A
 \simeq(B_\star A\star B_\star A)^\sharp
 \xrightarrow{\Delta_B^\sharp}(B_\star A)^\sharp.
\]
The bar counit becomes the unit; the coaugmentation becomes the
augmentation. This constructs a second associative algebra, with
the tensor in \eqref{cyjets:eq:dual-tensor}.

\begin{theorem}\label{cyjets:thm:bar-duality}
Assume that normalized associative bar and cobar are inverse
equivalences between the augmented algebras with positive
augmentation ideal in $(\Ccal,\star)$ and their connected
conilpotent coalgebras. Assume their normalized words have the
finite configuration bound above. Then $\mathcal K_\sharp$ is an
anti-equivalence to the corresponding augmented associative algebras
for $\star_\sharp$. Its inverse is
\[
 E\longmapsto\mathcal U B_{\star_\sharp}E.
\]
\end{theorem}
\begin{proof}
Bar gives a coalgebra equivalence by hypothesis, and
Proposition~\ref{cyjets:prop:operations} transposes coalgebras to
algebras for the transported tensor. To identify the inverse, let
$C$ be a connected conilpotent $\star$-coalgebra. On normalized
words,
\[
 \mathcal U B_{\star_\sharp}(C^\sharp)\simeq\Omega_\star C.
\]
A bar suspension dualizes to a cobar desuspension. Multiplication
dualizes to the coproduct of $C$, while deconcatenation dualizes
to concatenation. The full differentials are transposed with these
maps. At each $I$ there are only finitely many positive words,
so exact duality commutes with their finite extensions.
For $C=B_\star A$, the comparison back to $A$ is therefore
\[
 \mathcal U B_{\star_\sharp}(\mathcal K_\sharp A)
       \simeq\Omega_\star B_\star A\xrightarrow{\epsilon_A}A,
\]
the stipulated bar--cobar counit. The other comparison is the
coalgebra-side unit. Their triangle identities and coefficient
biduality prove the inverse-functor assertion.
\end{proof}

The theorem isolates coefficient transposition from the geometric
bar--cobar construction: its hypotheses must be proved for the
chosen collision convolution. They do not follow from Lie--commutative
chiral duality. Once supplied, the full bar coalgebra yields the
second associative algebra above, even for nonholonomic currents.
No quadratic approximation of the bar has been used.

At finite parameter precision the same argument uses
$D_X\otimes_kR_N$. Perfect modules remain the compact modules
over this associative ring, although $R_N$ is not regular.
Coefficient reduction is transported by \eqref{cyjets:eq:transport},
and the compatible inverse system is part of the dual datum.
The finite word bound does not exchange an ordinary pro-limit,
a parameter-adic limit and a Laurent expansion.

\section{The further field comparison}
\label{cyjets:sec:comparison}

Theorem~\ref{cyjets:thm:transpose} identifies the zero-potential
Hall cooperation with a field product on restricted jets.
Theorem~\ref{cyjets:thm:bar-duality} instead transposes a full
coalgebra to an algebra for a transported geometric tensor.
Their carriers and operations specify two different comparison
problems. Neither identifies a Hall charge split with a
level-changing field coproduct: a map sending a charge-one state
to two charge-one factors has output charges $(1,1)$, whereas a
Hall split has output charges adding to one.

The empty-configuration unit $K$ is also different from the local
vacuum coefficient $\omega_X$. An augmentation is a specified
algebra map to the tensor unit, and evaluation in a prepared state
is additional linear data. The coefficient duality constructs none
of these identifications by changing notation.

\subsection{The full bar contact marking}

Suppose a native field system has already supplied a connected full
collision coalgebra $C$ in the category above. Its contact maps retain
\[
 r(u^m;a,b)=\operatorname{pp}_u\bigl(u^mY(a,u)b\bigr),
                      \qquad m\in\Z,
\]
where $\operatorname{pp}_u$ takes the finite principal part at the
stated precision. The coalgebra also retains the higher collision
maps, spectators and normal derivatives. Form $U_*(C)=\Omega_\star C$.
Under the bar--cobar hypotheses of
Theorem~\ref{cyjets:thm:bar-duality}, the unit is the identity
at empty support and has the following formula on $c\in\bar C$:
\begin{equation}\label{cyjets:eq:bar-contact}
 \eta_C:C\xrightarrow{\sim}B_\star\Omega_\star C,
 \qquad
 \eta_C(c)=\sum_{r\geq1}(s\tau)^{\star r}
                           \bar\Delta^{r-1}c.
\end{equation}
Here $s$ has degree $-1$, $t$ has degree $1$, and $\tau(c)=tc$
is the universal twisting cochain into the cobar. At a configuration
$I$ the sum has $r\leq|I|$. With bar signs
$b_1(sa)=-s(da)$ and $b_2(sa,sb)=(-1)^{|a|}s(ab)$,
cobar splitting cancels adjacent bar multiplication. The remaining
terms give the chain identity for $\eta_C$; cutting the outer word
gives its coalgebra identity. The stipulated bar--cobar theorem
supplies its inverse equivalence.

This marking retains every native binary mode:
\begin{equation}\label{cyjets:eq:all-modes}
 a_{(n)}b=[u^{-1}]r(u^n;a,b),\qquad n\in\Z.
\end{equation}
Indeed multiplication by $u^n$ moves the coefficient of
$u^{-n-1}$ to $u^{-1}$. Negative multipliers are essential for
recovering regular modes. Transposition gives the operation-preserving
coefficient comparison
\[
 (B_\star\Omega_\star C)^\sharp
                 \xrightarrow{\eta_C^\sharp}C^\sharp.
\]
These operations are marked on the bar of the algebra realization.
They have not been identified with its bare singleton multiplication.
The higher native collision system was used to construct $C$;
equation \eqref{cyjets:eq:all-modes} alone does not construct it.

\subsection{Equal fields can carry different jet markings}

A further ambiguity remains even when all unmarked fields agree.
Take $R_2=k[\hbar]/(\hbar^2)$, with $\hbar$ of weight and cochain
degree zero, and
the charge-graded algebra
\[
 A=R_2[e_0,e_1,\ldots]/(e_0,e_1,\ldots)^3.
\]
Every $e_n$ has charge one, weight $n+1$ and cochain degree zero;
the differential is zero. There is no cutoff on $n$. Define
derivations by
\[
 Te_n=(n+1)e_{n+1},\qquad
 \delta_i e_n=e_{n-i+1}\quad(i\geq1),
\]
where $e_j=0$ for $j<0$. Then $\delta_1$ is the charge derivation,
the $\delta_i$ commute, and
$[\delta_i,T]=(i-1)\delta_{i-1}$, with $\delta_0=0$.
For each $\kappa\in k$, change the marking to
\begin{equation}\label{cyjets:eq:moment-marking}
 J_1^\kappa=\delta_1,\quad J_2^\kappa=\delta_2,\quad
 J_i^\kappa=\delta_i+
       \hbar\kappa\binom{i-1}{2}\delta_{i-2}\quad(i\geq3).
\end{equation}
Every operator is charge preserving and bounded for the weight
filtration. Only finitely many act nontrivially on any fixed state.
They commute with the differential and satisfy the same jet and
translation relations. The latter assertion follows from
\[
 \binom{i-1}{2}(i-3)=(i-1)\binom{i-2}{2}.
\]

For these markings define the contraction and fields
\begin{align*}
 C_1^\kappa(u)
 &=\sum_{p,q\geq0}(-1)^{p+1}\binom{p+q}{p}
       u^{-p-q-1}J_{p+1}^\kappa\otimes J_{q+1}^\kappa,\\
 Y^\kappa(a,u)b
 &=\mu(e^{uT}\otimes1)(1-\hbar C_1^\kappa(u))(a\otimes b),\\
 \mathcal S^\kappa(u)&=1+2\hbar C_1^\kappa(u).
\end{align*}
The contraction is finite on each pair of algebraic inputs.
The commuting derivations and their translation identities give
the field identities at this precision: in a three-input product
only the sum of the three pair contractions survives. This also
gives exchange and the block-fusion identities.

\begin{proposition}\label{cyjets:prop:marked-ambiguity}
All the displayed unmarked fields and exchanges are identical.
Distinct values of $\kappa$ remain inequivalent if $R_2$-linear
morphisms preserve the markings and weak equivalences are state quasi-isomorphisms
through nonpositive dg algebras whose marked operators are chain maps.
\end{proposition}
\begin{proof}
The difference $J_i^\kappa-\delta_i$ is divisible by $\hbar$.
Since the contraction already has an external factor $\hbar$,
its change is zero in $R_2$. Thus $Y^\kappa=Y^0$ and
$\mathcal S^\kappa=\mathcal S^0$. In particular all iterates and
Laurent-multiplier residues agree.

On charge one the standard operators satisfy
$\delta_3=\delta_2^2$. The modified marking has
\begin{equation}\label{cyjets:eq:moment-defect}
 J_3^\kappa-(J_2^\kappa)^2=\hbar\kappa\,\id.
\end{equation}
Chain compatibility makes this an operator on cohomology. A zigzag
of marked quasi-isomorphisms induces an invertible intertwiner $f$
on charge-one cohomology. Between parameters $\kappa,\lambda$,
equation \eqref{cyjets:eq:moment-defect} would force
$\hbar(\kappa-\lambda)f=0$. Composing with $f^{-1}$ contradicts
freeness of the charge-one $R_2$-module unless $\kappa=\lambda$.
\end{proof}

The assumption that the marked operators are chain maps is necessary
for this cohomological argument. Exact weight would also change the
admitted class: the correction in \eqref{cyjets:eq:moment-marking}
has weight $3-i$, whereas $\delta_i$ has weight $1-i$.
Filtered boundedness admits both. On the bottom state the difference
is already visible as $J_2e_0=0$ and $J_3^\kappa e_0=
\hbar\kappa e_0$, but the translation orbit remains infinite.

This ambiguity does not destroy a specified current realization.
For example, in charge $n$ let $p_m=\sum_{a=1}^nx_a^m$ and
$p_0=n$. The corresponding marked moments and their compensation are
\[
 P_m^\kappa=p_m+\hbar\kappa\binom m2p_{m-2},\qquad
 \widehat P_m=P_m^\kappa-
                   \hbar\kappa\binom m2P_{m-2}^\kappa=p_m,
\]
with the correction zero for $m<2$. Thus compensated moments give
the original action of $k[x_1,\ldots,x_n]^{\mathfrak S_n}$.
Any current presentation with this specified action retains its
localization at $\prod_{a<b}(x_a-x_b)^2$ and its rational pair
multipliers. Raw moments cannot replace that action: in charge one,
$P_2^\kappa-(P_1^\kappa)^2=\hbar\kappa$, while the polynomial
identity is $p_2=p_1^2$.

The bar contact comparison and the specified current markings are
therefore positive additional data. Recovering those markings
intrinsically from an unmarked field is a stronger problem. The
first-order equality above makes no assertion of compatible equality
at every higher parameter precision.

\begin{problem}\label{cyjets:prob:realization}
For an independently defined family of local fields, construct its
coherent collision coefficient system and compare its operation
maps with the transported full-bar algebra. Specify the coefficient
class, tensor, unit, translation and Laurent completion, and retain
the maps supported on collision diagonals. For critical Hall input,
construct the oriented coefficient map before applying the restricted
transpose; the zero-potential polynomial calculation alone does not
supply it.
\end{problem}

The next geometric input is the coefficient at a collision. The
derived relation directions can survive where their ordinary
equations vanish. Retaining those directions is necessary before
testing a comparison with the fields constructed here.
\endgroup
